%%%%%%%%%%%%%%%%%%%%%%%%%%%%%%%%%%%%%%%%%%%%%%%%%%%%%%%%%%%%%%%%%%%%%%%%%%%%%%%
% Janus: A Hierarchical Multi-Agent Framework with Structured Quality Assurance
% Technical Whitepaper — arXiv Preprint
%%%%%%%%%%%%%%%%%%%%%%%%%%%%%%%%%%%%%%%%%%%%%%%%%%%%%%%%%%%%%%%%%%%%%%%%%%%%%%%

\documentclass[11pt,a4paper]{article}

% ── Packages ────────────────────────────────────────────────────────────────
\usepackage[utf8]{inputenc}
\usepackage[T1]{fontenc}
\usepackage{newtxtext,newtxmath}       % Times-like font
\usepackage[margin=1in]{geometry}
\usepackage{graphicx}
\usepackage{listings}
\usepackage{xcolor}
\usepackage{booktabs}
\usepackage{hyperref}
\usepackage{enumitem}
\usepackage{float}
\usepackage{caption}
\usepackage{subcaption}
\usepackage{fancyhdr}
\usepackage{cite}
\usepackage{amsmath}
\usepackage{amssymb}

% ── Hyperref setup ──────────────────────────────────────────────────────────
\hypersetup{
    colorlinks=true,
    linkcolor=blue,
    citecolor=blue,
    urlcolor=blue,
}

% ── Code listing style ──────────────────────────────────────────────────────
\definecolor{codebg}{rgb}{0.95,0.95,0.95}
\lstdefinestyle{janus}{
    backgroundcolor=\color{codebg},
    basicstyle=\ttfamily\small,
    breaklines=true,
    frame=single,
    columns=fullflexible,
    keepspaces=true,
    language=Python,
    numbers=none,
    showstringspaces=false,
}
\lstset{style=janus}

% ── Page style ──────────────────────────────────────────────────────────────
\pagestyle{fancy}
\fancyhf{}
\fancyhead[L]{\small Janus: A Hierarchical Multi-Agent Framework}
\fancyhead[R]{\small \thepage}
\renewcommand{\headrulewidth}{0.4pt}

% ── Title ───────────────────────────────────────────────────────────────────
\title{\textbf{Janus: A Hierarchical Multi-Agent Framework\\[4pt]
        with Structured Quality Assurance}}

\author{%
    Bingsheng Wang\\
    Taiyuan University of Technology\\
    \texttt{https://github.com/isheng-eqi/janus}
}

\date{July 2026}

\begin{document}

\maketitle

\begin{abstract}
Janus is a hierarchical multi-agent framework for autonomous task execution
that organizes LLM-based agents into a four-role architecture inspired by human
organizational management rather than distributed systems design.
The framework introduces the \emph{Gatekeeper Tree} pattern:
a \textbf{Gatekeeper} (strategist with zero tools) sets direction and validates delivery;
a \textbf{Planner} (tactician with zero tools) decomposes directives into executable work packages;
\textbf{Workers} (executors with nine concrete tools) perform the actual work;
and a \textbf{Reviewer} (independent auditor with zero tools) inspects every
deliverable against structured acceptance criteria.
Janus introduces three core innovations: (1)~a five-tier graded review verdict
system (approved, approved-with-notes, minor revisions, major revisions,
rejected) coupled with a four-level defect severity classification (critical,
major, minor, suggestion), directly inspired by academic peer review and
manufacturing quality control; (2)~a full-chain intent-preservation mechanism
where the user's original request propagates unmodified through every layer,
serving as an immutable anchor for delivery validation; and (3)~a self-healing
closed loop with structured diagnosis, strategic reformulation, and result
merging across recovery attempts.
We evaluate Janus on complex multi-step software engineering tasks and
demonstrate that the hierarchical review pipeline catches errors that
flat single-agent and pipeline architectures miss, while the graded verdict
system prevents unnecessary retries on soft criteria and the intent loop
catches goal drift before results reach the user.
\end{abstract}

% =============================================================================
\section{Introduction}
% =============================================================================

The shift from monolithic large language models (LLMs) to LLM-based autonomous
agents represents one of the most consequential developments in applied AI.
Rather than treating a model as a question-answering oracle, agent frameworks
equip LLMs with tools, memory, and decision-making loops, enabling them to read
files, execute commands, search the web, and iteratively refine their output.
This paradigm has produced frameworks such as LangGraph~\cite{langgraph},
AutoGen~\cite{autogen}, CrewAI~\cite{crewai}, and MetaGPT~\cite{metagpt},
each offering a different approach to orchestrating multiple LLM calls toward
a shared goal.

\subsection{The Flatness Problem}

Despite this rapid progress, a structural weakness runs through nearly all
existing agent frameworks: they are \emph{architecturally flat}.
In a typical LangGraph or AutoGen pipeline, every agent node shares the same
tool surface (the ability to read, write, and execute), the same
responsibility scope (both plan and execute), and---critically---the same
information horizon (full access to all intermediate results).
While this flatness makes frameworks easy to reason about, it introduces
three interrelated failure modes:

\begin{enumerate}[leftmargin=*,itemsep=2pt]
    \item \textbf{Hallucination cascades.}
    When every agent can read every intermediate output, an early hallucination
    propagates unchecked through subsequent steps. A worker that misidentifies
    a target directory contaminates every downstream worker with the wrong
    context, and no independent agent exists to catch the drift.

    \item \textbf{Quality collapse under scale.}
    Existing frameworks lack \emph{structured} quality gates. Review, when it
    exists at all, is a binary pass/fail decision made by the same LLM that
    produced the output---a violation of the separation-of-concerns principle
    known to every human organization. The result is that complex multi-step
    tasks degrade into correctness-by-accident.

    \item \textbf{Intent drift.}
    As a user request passes through multiple LLM-to-LLM handoffs---from
    orchestrator to decomposer to worker---each layer introduces a small
    interpretation bias. By the time output reaches the user, the original
    intent may have silently mutated. No existing framework systematically
    guards against this.
\end{enumerate}

\subsection{The Janus Insight}

Human organizations have spent millennia solving exactly these problems.
Armies separate strategic command from tactical execution.
Governments establish independent inspector generals to audit administrative
agencies. Manufacturers deploy three-stage quality inspection (incoming,
in-process, outgoing). Academic publishers grade peer review outcomes into
accept/minor-revisions/major-revisions/reject.
Courts apply different standards of review depending on what is being examined.

Janus's core insight is that \emph{LLM agent architecture should mirror human
organizational design, not distributed systems design}.
Where LangGraph and AutoGen model their architectures on workflow graphs and
message-passing topologies, Janus models its architecture on the
separation-of-concerns, accountability hierarchies, and structured quality
assurance that human institutions have refined over centuries.

This paper presents five principal contributions:

\begin{itemize}[leftmargin=*,itemsep=2pt]
    \item \textbf{The Gatekeeper Tree architecture:} a four-role hierarchy
    (Gatekeeper, Planner, Worker, Reviewer) with hard tool-access constraints
    per role and layered information horizons. Gatekeeper, Planner, and
    Reviewer each have \emph{zero tools}---only Workers may read files, write
    files, or execute commands. This forces each role to operate at its
    designated level of abstraction.

    \item \textbf{A five-tier graded review system:} inspired by academic peer
    review (accept/minor-revisions/major-revisions/reject) and manufacturing
    quality control (critical/major/minor/suggestion defect severity),
    replacing binary pass/fail with nuanced verdicts and tiered retry
    strategies that prevent unnecessary re-execution.

    \item \textbf{Full-chain intent preservation:} the \texttt{user\_goal}
    field propagates the user's original, unmodified request through every
    layer---Gatekeeper, Planner, Worker, Reviewer, and back---serving as an
    immutable anchor against which final delivery is validated before the
    user ever sees it.

    \item \textbf{A self-healing recovery loop:} when tasks fail, the
    Gatekeeper performs structured diagnosis (asking the LLM \emph{why} tasks
    failed), strategic reformulation (creating a new directive targeting only
    the failed aspects), re-execution, and result merging across attempts.

    \item \textbf{An empirical comparison} with existing frameworks
    (LangGraph, AutoGen, CrewAI, MetaGPT) across architectural, quality,
    safety, and observability dimensions, plus an end-to-end case study
    demonstrating the full pipeline from user input to validated output.
\end{itemize}

% =============================================================================
\section{Design Philosophy}
% =============================================================================

This section articulates \emph{why} Janus is structured the way it is.
We establish four design principles derived from human organizational patterns,
then show how each principle maps to a concrete architectural constraint.

\subsection{Why Hierarchy?}

An LLM's context window is both its greatest asset and its tightest constraint.
Every token loaded into context consumes attention budget and risks diluting
the signal-to-noise ratio. A flat architecture---where every agent sees every
intermediate result---is optimal for information completeness but catastrophic
for information \emph{relevance}. The agent that must decide whether a task
succeeded should not be the same agent that reads every line of tool-call
output from every worker.

Human institutions converged on hierarchy for exactly this reason. A general
does not read individual soldiers' after-action reports; a CEO does not
review every line of code committed by every engineer. Each layer compresses
and interprets the layer below it, adding \emph{judgment} while removing
\emph{noise}.

\subsection{Four Design Principles}

\subsubsection*{Principle 1: Separation of Concerns}
\label{sec:sep-concerns}

No single agent should both produce work and judge its quality.
In Janus, the Gatekeeper sets direction but never executes; the Worker executes
but never self-reviews; the Reviewer audits but never modifies.
We enforce this at the \emph{tool level}: Gatekeeper, Planner, and Reviewer are
initialized with zero tools---they cannot read files, write files, or run
commands. Only Workers have tool access. This is not a suggestion; it is a
hard architectural constraint enforced by the \texttt{ToolRegistry}.

This principle maps directly to the military Staff/Line split (planners do not
command troops) and the government Inspector General model (auditors are
independent of the agencies they audit).

\subsubsection*{Principle 2: Trust Nothing}
\label{sec:trust-nothing}

Every Worker output must pass independent review before the Gatekeeper
presents it to the user. The Reviewer is a separate LLM instance with a
dedicated system prompt, no shared context with the Worker, and strict
instructions to inspect evidence, not assume it. We model this on
manufacturing's three-line quality defense: incoming inspection (TaskSpec
validation), in-process inspection (desk-check pre-screening), and outgoing
inspection (full Reviewer audit).

\subsubsection*{Principle 3: Recursive Decomposition}
\label{sec:recursive-decomp}

Complex goals are recursively broken into sub-goals until each sub-goal is
small enough for a single Worker to handle in one pass. The decomposition tree
has a hard depth limit (configurable, default 3) to prevent infinite recursion.
Workers can self-decompose when a task proves too complex mid-execution,
requesting decomposition and receiving sub-task results before resuming.

This principle maps to the OKR (Objectives and Key Results) cascade from
corporate management: each level of decomposition must explicitly state
\emph{why} the sub-task matters (the \texttt{intent} field), not just
\emph{what} to do.

\subsubsection*{Principle 4: Human-in-the-Loop by Default}
\label{sec:human-in-loop}

The Gatekeeper's \texttt{\_report\_to\_user} method presents only
architecture-level information: how many tasks passed, how many failed,
a one-line summary per task. The user never sees raw Worker output, tool-call
logs, or implementation details unless explicitly requested (via \texttt{--verbose} mode).
This respects the user's cognitive bandwidth while preserving full
traceability for debugging. The delivery validation step
(\texttt{\_validate\_delivery}) runs before user-facing output, catching
goal drift before it reaches the user.

\subsection{The Human-Organization Design Pipeline}

Janus formalizes a four-step process for designing new features:

\begin{enumerate}[leftmargin=*,itemsep=2pt]
    \item \textbf{Find the human equivalent.} Ask: ``What does this problem
    look like in a human organization?'' Map to military, government,
    corporate, judicial, manufacturing, or academic practice.
    \item \textbf{Extract the core mechanism.} Strip away institutional
    scaffolding (we do not need military tribunals or HR departments).
    Keep only what directly improves decision quality or execution
    reliability.
    \item \textbf{Map to Janus roles and protocols.} Translate the human
    concept into a concrete data structure or method. For example, Commander's
    Intent $\rightarrow$ \texttt{TaskSpec.intent}; Inspector General
    $\rightarrow$ \texttt{Reviewer}; after-action report $\rightarrow$
    \texttt{ExecutionReport}.
    \item \textbf{Implement minimally.} Add only what the current system
    demonstrably needs. Defer multi-planner parallelism, resource arbitration,
    and cross-worker negotiation until those problems actually manifest.
\end{enumerate}

This is a \emph{design methodology}, not a metaphor. Every structural
decision in Janus---from the five-tier review verdict to the context-discipline
prompt---can be traced to a specific human practice that was studied, extracted,
and mapped.

% =============================================================================
\section{Janus Architecture}
% =============================================================================

This section provides a detailed technical description of Janus's four-role
architecture, the information flow between roles, and the protocol data
structures that carry state across the system. We describe what each layer
sees, what it produces, and what it is deliberately prevented from seeing.

\subsection{System Overview}

Figure~\ref{fig:architecture} shows the complete Janus architecture.
The user interacts with a REPL loop via \texttt{Session}, which maintains
conversation history and delegates all decisions to the Gatekeeper.
The Gatekeeper routes simple chat directly to a lightweight LLM call, and
complex tasks through the full Planner $\rightarrow$ Worker $\rightarrow$
Reviewer pipeline.

\begin{figure}[H]
\centering
\begin{verbatim}
                         User
                          |
                       Session (history + pass-through)
                          |
                    ┌─────Gatekeeper──────┐
                    │  Strategic layer    │
                    │  ZERO tools         │
                    │  Model: reasoning   │
                    └────────┬────────────┘
                             | Directive
                    ┌────────Planner──────┐
                    │  Tactical layer     │
                    │  ZERO tools         │
                    │  Model: fast        │
                    └──┬──────┬──────┬────┘
                       |      |      |
                    Worker0 Worker1 WorkerN
                    (tools) (tools) (tools)
                       |      |      |
                    ┌──┴──────┴──────┴────┐
                    │     Reviewer        │
                    │  Independent audit  │
                    │  ZERO tools         │
                    └─────────────────────┘
\end{verbatim}
\caption{Janus four-role architecture. Only Workers have tool access;
Gatekeeper, Planner, and Reviewer operate through pure LLM reasoning.}
\label{fig:architecture}
\end{figure}

\subsection{Gatekeeper: The Strategic Decision Layer}

The Gatekeeper (\texttt{core/gatekeeper.py}, 973 lines) is the sole entry
point for all user input. Its design is governed by a single hard constraint:
\emph{zero tools}. The Gatekeeper cannot read files, write files, execute
commands, or search the web. It reasons purely through LLM calls.

\subsubsection*{Decision Routing}

The \texttt{\_decide} method classifies every user message as either
\texttt{chat} (simple conversation) or \texttt{task} (needs Planner
decomposition). The decision is a structured JSON call:

\begin{lstlisting}
{
    "action": "chat" | "task",
    "reason": "brief explanation"
}
\end{lstlisting}

API failures default to \texttt{task} (fail-open: it is better to attempt
decomposition and fail gracefully than to silently treat a complex request as
chat). The decision prompt includes a \texttt{context\_discipline\_prompt}
shared across all strategic layers, which instructs the LLM to maintain only
architecture-level information and summarize Worker results to one line.

\subsubsection*{Directive Formulation}

When the action is \texttt{task}, the Gatekeeper calls
\texttt{\_formulate\_directive} to translate the user's goal into a
\texttt{Directive} object (Section~\ref{sec:protocol}) with three fields:
\texttt{intent} (strategic direction), \texttt{constraints} (hard must-not-violate
rules), and \texttt{priority} (speed, quality, or balanced). The
\texttt{user\_goal} field preserves the user's original input verbatim---it is
never modified by any LLM call---and serves as the system's immutable anchor.

A critical detail: the \texttt{user\_goal} is separate from \texttt{goal}.
During recovery loops, \texttt{goal} may be augmented with diagnostic
annotations (``[RECOVERY ATTEMPT 1/2 --- only address failed aspects]''),
but \texttt{user\_goal} remains pristine. When delivery validation runs
(Section~\ref{sec:intent-loop}), it compares output against
\texttt{user\_goal}, not against the possibly-augmented \texttt{goal}.

\subsubsection*{Recovery Loop}

If the Planner returns an \texttt{ExecutionReport} with failed tasks, the
Gatekeeper enters a structured recovery loop (up to 2 attempts):

\begin{enumerate}[leftmargin=*,itemsep=2pt]
    \item \textbf{Diagnose} (\texttt{\_diagnose\_failures}): ask the LLM
    \emph{why} tasks failed, using structured failure data from the report
    including acceptance criteria and review issues.
    \item \textbf{Reformulate} (\texttt{\_reformulate\_for\_recovery}): create
    a new \texttt{Directive} targeting \emph{only} the failed aspects, with an
    explicitly different strategy based on the diagnosis.
    \item \textbf{Re-execute}: pass the new directive to the Planner.
    \item \textbf{Merge} (\texttt{\_merge\_reports}): combine successful
    results from the original attempt with new results from recovery.
\end{enumerate}

This corresponds to the corporate Performance Improvement Plan (PIP) pattern:
not just ``redo it,'' but ``diagnose why it failed, specify what must change,
and try a different approach.''

\subsubsection*{Delivery Validation}

Before formatting the final user-facing output, the Gatekeeper runs
\texttt{\_validate\_delivery}, a single LLM call that asks: ``Does the output
actually answer the user's original request?'' If the answer is no, the
Gatekeeper triggers an extra recovery cycle. This is the final safety net
for intent drift (Section~\ref{sec:intent-loop}).

\subsection{Planner: The Tactical Execution Layer}

The Planner (\texttt{core/planner.py}, 1,155 lines) receives a
\texttt{Directive} from the Gatekeeper and produces an
\texttt{ExecutionReport}. Like the Gatekeeper, the Planner has \emph{zero
tools}---it cannot read, write, or execute. It only plans, dispatches, tracks,
and summarizes.

\subsubsection*{Tactical Decomposition}

The \texttt{\_plan} method calls the LLM to break the directive into an array
of \texttt{TaskSpec} objects. Each \texttt{TaskSpec} contains:

\begin{itemize}[leftmargin=*,itemsep=2pt]
    \item \texttt{task\_id}: unique identifier.
    \item \texttt{description}: concrete, actionable instruction.
    \item \texttt{acceptance\_criteria}: each criterion prefixed with
    \texttt{[HARD]} (must-have, zero tolerance) or \texttt{[SOFT]}
    (nice-to-have, minor deviations acceptable). Unmarked criteria default
    to \texttt{[HARD]} for backward compatibility.
    \item \texttt{context}: scoped background information (not full history).
    \item \texttt{intent}: why this task matters in the bigger picture.
    \item \texttt{goal}, \texttt{user\_goal}, \texttt{constraints}: propagated
    from the directive.
    \item \texttt{depth}: current recursion depth (1 = root task).
\end{itemize}

The \texttt{[HARD]/[SOFT]} classification is a direct mapping of judicial
review standards: \texttt{[HARD]} criteria receive \emph{de novo} scrutiny
(zero tolerance), while \texttt{[SOFT]} criteria receive a ``clearly
erroneous'' standard (only blatant failures block). This prevents the system
from wasting retry cycles on cosmetic issues like variable naming style when
the core functionality is correct.

\subsubsection*{Priority-Driven Retry Budgets}

The Planner maps the directive's \texttt{priority} field to a concrete retry
budget:

\begin{center}
\begin{tabular}{lc}
\toprule
\textbf{Priority} & \textbf{Max Retries} \\
\midrule
speed / urgent & 0 (fail fast, no retries) \\
balanced / normal & 2 (up to 3 total attempts) \\
quality & 3 (up to 4 total attempts) \\
\bottomrule
\end{tabular}
\end{center}

This allows users to trade execution thoroughness against latency.

\subsubsection*{Dispatch with Review}

The \texttt{\_dispatch\_with\_review} method implements a graded retry
strategy keyed to the Reviewer's verdict (Section~\ref{sec:reviewer}):

\begin{center}
\begin{tabular}{ll}
\toprule
\textbf{Verdict} & \textbf{Action} \\
\midrule
APPROVED / APPROVED\_WITH\_NOTES & Accept immediately \\
MINOR\_REVISIONS & Retry once with feedback, auto-accept after \\
MAJOR\_REVISIONS & Retry up to 2$\times$ with full re-review \\
REJECTED & Retry up to 2$\times$, then fail \\
\bottomrule
\end{tabular}
\end{center}

Additionally, after the Reviewer returns its verdict, the
\texttt{\_adjust\_verdict\_for\_criteria} method applies a safety net:
if all blocking issues are \texttt{MINOR} severity but the verdict is
\texttt{MAJOR\_REVISIONS} or \texttt{REJECTED}, the verdict is demoted to
\texttt{MINOR\_REVISIONS}---the Reviewer LLM was too harsh.

\subsubsection*{Desk Check}

Before sending any Worker result to the Reviewer, the Planner runs a
zero-cost \texttt{\_desk\_check}: a set of heuristic rules that flag
suspicious results without an LLM call. A Worker result that is empty,
appears to be just a file path, or lists artifacts but has a generic summary
(``done'' / ``ok'') triggers a warning injected into the review context.
This is modeled on academic ``desk reject'': the editor checks paper quality
\emph{before} sending to reviewers.

\subsection{Worker: The Tool-Execution Layer}
\label{sec:worker}

The Worker (\texttt{core/worker.py}, 1,363 lines) is the only role with
access to tools. Each Worker runs an LLM-driven loop: receive a
\texttt{TaskSpec}, call tools to accomplish the task, and return a
\texttt{TaskResult} as structured JSON.

\subsubsection*{Tool Registry and 9 Tools}

The \texttt{ToolRegistry} maps tool names to Python callables and generates
OpenAI-compatible function-calling schemas. A parameter alias system
(\texttt{\_PARAM\_ALIASES}) maps LLM-preferred parameter names (e.g.,
\texttt{command} $\rightarrow$ \texttt{cmd}, \texttt{filename} $\rightarrow$
\texttt{path}) to canonical function signatures, filtered through
\texttt{inspect.signature} to prevent unexpected arguments from reaching
underlying functions.

The default registry (\texttt{create\_default\_registry}) registers nine tools:

\begin{center}
\begin{tabular}{lll}
\toprule
\textbf{Tool} & \textbf{Function} & \textbf{Limit} \\
\midrule
\texttt{read\_file} & Read file content with offset/limit paging & 50,000 chars/call \\
\texttt{write\_file} & Write content, auto-create parent dirs & None \\
\texttt{terminal} & \texttt{subprocess} shell execution & 60s timeout \\
\texttt{web\_search} & Web search (placeholder) & --- \\
\texttt{web\_extract} & HTTP GET + HTML strip & 3,000 chars/URL, max 5 URLs \\
\texttt{search\_files} & Glob + content search in CWD & 50 results \\
\texttt{patch} & Find-and-replace in file & First occurrence only \\
\texttt{execute\_code} & \texttt{exec()} Python in sandboxed namespace & Restricted builtins \\
\texttt{browser\_navigate} & Browser navigation (placeholder) & --- \\
\bottomrule
\end{tabular}
\end{center}

All tool crashes are caught; the Worker loop never dies from a misbehaving tool.

\subsubsection*{Self-Decomposition}

When a Worker determines mid-execution that a task is too complex to complete
in one pass, the LLM can return \texttt{status = "needs\_decomposition"} with
a \texttt{decomposition\_request} containing sub-tasks. The Worker recursively
executes each sub-task (bounded by \texttt{max\_depth}, default 3), feeds
sub-results back as context, and resumes the original task. Only one level of
self-decomposition is permitted per task; a second
\texttt{NEEDS\_DECOMPOSITION} after resume produces a \texttt{FAILURE}.

Each sub-task result passes through \texttt{\_review\_sub\_result}, which
applies the same graded retry logic as the Planner-level review loop.

\subsection{Reviewer: The Independent Audit Layer}
\label{sec:reviewer}

The Reviewer (\texttt{core/reviewer.py}, 623 lines) is a standalone LLM agent
with \emph{zero tools}. It receives a \texttt{TaskSpec} and a
\texttt{TaskResult}, then evaluates whether the result meets every acceptance
criterion. The Reviewer is called by the Planner after every Worker execution
(including sub-Worker executions during self-decomposition).

\subsubsection*{Artifact Pre-Loading}

Before calling the LLM, the Reviewer reads artifact files declared in the
Worker's \texttt{artifacts} list and embeds their contents directly into the
review prompt (up to 1,000 characters per file, 8,000 total). This allows the
LLM to inspect what was \emph{actually} written to disk, not just what the
Worker \emph{claimed} to write. Access control limits file reading to only
the paths the Worker itself declared as artifacts, preventing path traversal.

\subsubsection*{Five-Tier Review Verdict}

The review outcome uses a five-tier verdict system defined in the
\texttt{ReviewVerdict} enum:

\begin{center}
\begin{tabular}{lll}
\toprule
\textbf{Verdict} & \textbf{Meaning} & \textbf{Behavior} \\
\midrule
\texttt{APPROVED} & All criteria perfectly met & Pass immediately \\
\texttt{APPROVED\_WITH\_NOTES} & All HARD criteria met, only SOFT suggestions & Record notes, pass \\
\texttt{MINOR\_REVISIONS} & Small issues, especially SOFT-only & Retry once, auto-accept \\
\texttt{MAJOR\_REVISIONS} & One or more HARD criteria not met & Retry up to 2$\times$, re-review \\
\texttt{REJECTED} & Core HARD criteria violated & Retry up to 2$\times$, then fail \\
\bottomrule
\end{tabular}
\end{center}

This maps directly to academic peer review (accept / minor revisions / major
revisions / reject) and judicial appellate standards.

\subsubsection*{Four-Level Defect Severity}

Each issue found is tagged with a \texttt{Severity} level:

\begin{center}
\begin{tabular}{lll}
\toprule
\textbf{Severity} & \textbf{Meaning} & \textbf{Effect on Retry} \\
\midrule
\texttt{CRITICAL} & \ding{72} Fatal---output completely unusable & Always triggers retry \\
\texttt{MAJOR} & \ding{110} Serious---core requirement not met & Triggers retry \\
\texttt{MINOR} & \ding{108} Moderate---partial deviation, output usable & Retry once, then accept \\
\texttt{SUGGESTION} & \ding{46} Optimization idea & Never blocks \\
\bottomrule
\end{tabular}
\end{center}

This maps to manufacturing defect classification: critical (stop the line),
major (quarantine batch), minor (record for rework), acceptable deviation
(concession).

\subsection{Inter-Layer Communication Protocol}
\label{sec:protocol}

Communication between layers is governed by four principal data structures
defined in \texttt{core/protocol.py} (256 lines):

\begin{itemize}[leftmargin=*,itemsep=2pt]
    \item \textbf{Directive} (Gatekeeper $\rightarrow$ Planner):
    \texttt{goal} (user's original or augmented goal),
    \texttt{user\_goal} (pristine, immutable),
    \texttt{intent} (strategic direction),
    \texttt{constraints} (hard must-not-violate rules),
    \texttt{priority} (speed/quality/balanced),
    \texttt{context} (multi-turn conversation history).

    \item \textbf{TaskSpec} (Planner $\rightarrow$ Worker):
    \texttt{task\_id}, \texttt{description}, \texttt{acceptance\_criteria}
    (with \texttt{[HARD]/[SOFT]} markers), \texttt{context}, \texttt{intent},
    \texttt{goal}, \texttt{user\_goal}, \texttt{constraints}, \texttt{depth}.

    \item \textbf{TaskResult} (Worker $\rightarrow$ Planner $\rightarrow$
    Reviewer):
    \texttt{status} (SUCCESS/FAILURE/NEEDS\_DECOMPOSITION),
    \texttt{summary}, \texttt{result}, \texttt{artifacts}, \texttt{confidence}
    (HIGH/MEDIUM/LOW), \texttt{decomposition\_request} (optional).

    \item \textbf{ExecutionReport} (Planner $\rightarrow$ Gatekeeper):
    \texttt{status} (completed/partial/failed),
    \texttt{total\_tasks}, \texttt{passed}, \texttt{failed},
    \texttt{summary}, \texttt{details} (per-task), \texttt{failed\_details},
    \texttt{failed\_tasks} (structured failure data for diagnosis).
\end{itemize}

The \texttt{user\_goal} field is the system's anchor: it propagates
\texttt{Session $\rightarrow$ Gatekeeper $\rightarrow$ Directive $\rightarrow$
Planner $\rightarrow$ TaskSpec $\rightarrow$ Worker $\rightarrow$ Reviewer
$\rightarrow$ back to Gatekeeper} without modification. The Gatekeeper's
delivery validation compares output against this field, not against any
LLM-interpreted version.

\subsection{Memory and State Management}

Janus uses three coordinated state-keeping mechanisms:

\begin{itemize}[leftmargin=*,itemsep=2pt]
    \item \textbf{Session} (\texttt{core/session.py}, 125 lines): a thin
    pass-through that records user/assistant message pairs (up to
    \texttt{max\_history = 100} turn-pairs) and formats recent history as
    context for the Gatekeeper's decision-making. Session makes no decisions
    of its own---it is purely a history recorder.

    \item \textbf{TaskManager} (\texttt{core/task\_manager.py}, 220 lines):
    a finite state machine tracking each task through
    \texttt{PENDING $\rightarrow$ RUNNING $\rightarrow$ COMPLETED/FAILED}.
    Transitions are validated against an allowed-transition map; illegal
    transitions raise \texttt{ValueError}. The TaskManager is reset at the
    start of each \texttt{Planner.execute()} call.

    \item \textbf{Console} (\texttt{core/console.py}, 466 lines): a passive
    observer that formats and displays events (phase transitions, task starts,
    tool calls, review results) with Chinese text, emoji, and ANSI color.
    Three output modes are supported: \texttt{default} (L0+L1+L2),
    \texttt{verbose} (includes full tool parameters), and \texttt{quiet}
    (only final summary).
\end{itemize}

% =============================================================================
\section{Key Mechanisms}
% =============================================================================

This section describes five mechanisms that collectively distinguish Janus
from flat agent architectures. Each mechanism addresses a specific failure
mode identified in the Introduction.

\subsection{Task Decomposition: The Gatekeeper Tree}

Janus's decomposition strategy is a \emph{recursive tree}, not a flat list.
The Gatekeeper formulates a Directive; the Planner decomposes it into
\texttt{TaskSpec[]}; Workers may further self-decompose mid-execution.
At every decomposition step, the \texttt{depth} field increments. When
\texttt{depth $\geq$ max\_depth} (default 3), self-decomposition is
disallowed and the task fails with an explicit depth-limit error.

This prevents the ``infinite recursion'' failure mode where a Worker
repeatedly declares a task too complex and spawns ever-smaller sub-tasks
without converging. The depth limit is uniform across all layers because the
Planner injects \texttt{self.\_max\_depth} into every Worker it creates.

Each \texttt{TaskSpec} carries an \texttt{intent} field---modeled on military
Commander's Intent---that explains \emph{why} the sub-task matters. This
allows Workers to make autonomous decisions when encountering unexpected
conditions (missing files, changed APIs), because they know the task's purpose
in the larger context.

\subsection{The Quality Audit Pipeline}
\label{sec:quality-audit}

Janus's quality assurance chain operates at three levels:

\begin{enumerate}[leftmargin=*,itemsep=2pt]
    \item \textbf{TaskSpec validation (IQC --- Incoming Quality Control).}
    Before any Worker is dispatched, the \texttt{TaskSpec.validate()} method
    checks that required fields (\texttt{task\_id}, \texttt{description}) are
    non-empty. Invalid specs are skipped with a warning.

    \item \textbf{Desk check pre-screening (IPQC --- In-Process Quality Control).}
    After Worker execution but before Reviewer audit, the
    \texttt{\_desk\_check} method applies zero-cost heuristics: is the result
    empty? Does it appear to be just a file path? Are artifacts present but
    the summary is generic? Warnings are injected into the review context
    for heightened scrutiny.

    \item \textbf{Full Reviewer audit (OQC --- Outgoing Quality Control).}
    The Reviewer LLM inspects the Worker result against every
    \texttt{[HARD]} and \texttt{[SOFT]} criterion, with artifact contents
    pre-loaded into the prompt for evidence-based evaluation.
\end{enumerate}

This three-tier approach mirrors manufacturing's quality defense system, where
defects caught at incoming inspection cost 1$\times$ to fix, in-process defects
cost 10$\times$, and post-delivery defects cost 100$\times$.

\subsection{The Intent Preservation Loop}
\label{sec:intent-loop}

Intent drift---the gradual mutation of the user's original request as it
passes through multiple LLM calls---is a well-documented failure mode in
multi-agent systems. Janus addresses this through a dual mechanism:

\begin{enumerate}[leftmargin=*,itemsep=2pt]
    \item \textbf{The \texttt{user\_goal} anchor.} At the moment the user's
    input enters the Gatekeeper, it is stored as \texttt{user\_goal} in the
    \texttt{Directive}. This field is \emph{never modified by any LLM call}.
    It propagates through every \texttt{TaskSpec} to every Worker, and appears
    in the Worker's system prompt as ``User's Original Request (do not
    modify).'' The Reviewer prompt also includes \texttt{user\_goal} with the
    explicit instruction: ``Confirm the output is what the user actually
    wanted, not just what satisfies the acceptance criteria.''

    \item \textbf{Delivery validation.} Before formatting the final user-facing
    output, \texttt{\_validate\_delivery} asks the LLM: ``User's original
    request: \texttt{\{user\_goal\}}. We produced: \texttt{\{report\_summary\}}.
    Does this output actually answer the user's request?'' If validation fails,
    the Gatekeeper triggers an extra recovery cycle.
\end{enumerate}

This mechanism was added after real-world incidents where the Planner's
unconditional working-directory injection caused Workers to analyze the wrong
project. The \texttt{user\_goal} anchor and pre-delivery validation now catch
such drift before the user ever sees incorrect output.

\subsection{The Self-Healing Closed Loop}

Janus's self-healing combines three feedback cycles at different architectural
levels:

\begin{enumerate}[leftmargin=*,itemsep=2pt]
    \item \textbf{Worker-level retry (Planner $\rightarrow$ Reviewer
    $\rightarrow$ Worker).} When the Reviewer returns
    \texttt{MINOR\_REVISIONS}, \texttt{MAJOR\_REVISIONS}, or
    \texttt{REJECTED}, the Planner constructs a retry \texttt{TaskSpec} with
    review feedback injected into the context. The retry prompt explicitly
    instructs the Worker: ``For each issue you fix, include proof:
    \texttt{\checkmark\ Fixed [issue]: what you changed and why it now
    satisfies the requirement.}''

    \item \textbf{Planner-level recovery (Gatekeeper $\rightarrow$ Planner).}
    When the \texttt{ExecutionReport} shows failed tasks, the Gatekeeper
    diagnoses \emph{why}, reformulates a new directive targeting only the
    failed aspects (preserving successes from the original attempt), and
    re-dispatches.

    \item \textbf{Delivery-level correction (Gatekeeper $\rightarrow$
    Planner).} When \texttt{\_validate\_delivery} detects intent drift, the
    Gatekeeper appends the deviation reason to the goal and triggers one
    final Planner execution.
\end{enumerate}

Results from multiple attempts are merged by \texttt{\_merge\_reports}, which
preserves successful tasks from all attempts and presents cumulative progress
to the user.

\subsection{Desk Check Pre-Screening}

The \texttt{\_desk\_check} method in \texttt{core/planner.py} (lines 944--998)
applies three zero-cost heuristics to every Worker result before it reaches
the Reviewer:

\begin{itemize}[leftmargin=*,itemsep=2pt]
    \item \textbf{Empty result:} If \texttt{result.result} is empty after
    stripping whitespace, flag ``Worker result is empty --- may be incomplete
    or failed silently.''
    \item \textbf{Path-only result:} If the result is a single short string
    containing a path separator (\texttt{/} or \texttt{\textbackslash}) or
    ending in a file extension, flag ``Worker result appears to be just a
    file path rather than substantive output.''
    \item \textbf{Generic summary with artifacts:} If artifacts are present
    but the summary is in a set of generic phrases (\texttt{\{``done,''
    ``ok,'' ``success,'' ``completed''\}}), flag ``Worker listed N artifacts
    but summary is generic.''
\end{itemize}

These heuristics cost zero LLM calls and catch the most common failure
patterns---empty output, file-path regurgitation, and missing
descriptions---before the Reviewer expends tokens on a full audit.

% =============================================================================
\section{Comparison with Existing Frameworks}
% =============================================================================

Table~\ref{tab:comparison} compares Janus with four major agent frameworks
across architectural, quality, safety, and observability dimensions.
We evaluate each framework as described in its public documentation and
source code at the time of writing.

\begin{table}[H]
\centering
\caption{Architectural comparison of Janus with existing agent frameworks.}
\label{tab:comparison}
\small
\begin{tabular}{p{2.3cm}p{2.3cm}p{2.3cm}p{2.3cm}p{2.3cm}p{2.3cm}}
\toprule
\textbf{Dimension} & \textbf{Janus} & \textbf{LangGraph} & \textbf{AutoGen} & \textbf{CrewAI} & \textbf{MetaGPT} \\
\midrule
\textbf{Architecture pattern}
    & Hierarchical tree (4 fixed roles)
    & Directed graph (nodes + edges)
    & Conversational (agent-to-agent messages)
    & Role-based team (flat)
    & Role-based pipeline (SOP-driven) \\
\midrule
\textbf{Tool access control}
    & Hard: only Workers have tools
    & Any node can have tools
    & Any agent can have tools
    & Any agent can have tools
    & Any agent can have tools \\
\midrule
\textbf{Quality review}
    & Independent Reviewer, 5-tier verdict, 4-level severity
    & Manual (user writes check nodes)
    & None built-in
    & None built-in
    & Code review role (binary) \\
\midrule
\textbf{Self-healing}
    & 3-layer closed loop (Worker $\rightarrow$ Planner $\rightarrow$ Gatekeeper)
    & Conditional edges (explicit retry paths)
    & Manual retry via group chat
    & None built-in
    & None built-in \\
\midrule
\textbf{Intent preservation}
    & \texttt{user\_goal} anchor + delivery validation
    & None
    & None
    & None
    & None \\
\midrule
\textbf{Information horizon}
    & Layered: each role sees only its layer's data
    & Full: all nodes share state graph
    & Full: all agents in group chat
    & Full: shared task context
    & Process-shared workspace \\
\midrule
\textbf{Design philosophy}
    & Human organizational management
    & Workflow graphs (Airflow/DAG)
    & Multi-agent conversation
    & Role-based team simulation
    & Software company SOP simulation \\
\midrule
\textbf{Extensibility}
    & Add roles only when demonstrably needed
    & Add nodes/edges arbitrarily
    & Add agents via register
    & Add agents via Crew definition
    & Add agents via Role classes \\
\midrule
\textbf{Observability}
    & 3-mode Console (default/verbose/quiet) with emoji
    & LangSmith tracing
    & Logging + callbacks
    & Logging
    & Logging + workspace \\
\bottomrule
\end{tabular}
\end{table}

\subsection{Key Differentiators}

\textbf{Where LangGraph excels:} LangGraph's directed-graph model is the most
general of the compared frameworks---it can express any control flow topology,
including Janus's tree structure. However, this generality comes at the cost
of built-in quality mechanisms: every retry path, every review node, and every
information filter must be hand-crafted by the developer. Janus provides these
as first-class architectural primitives with zero-configuration defaults.

\textbf{Where AutoGen excels:} AutoGen's multi-agent conversation model enables
rich, unstructured collaboration between agents that can debate, negotiate, and
iteratively refine output. For tasks requiring creative collaboration, this
flexibility is valuable. Janus's hierarchical model prioritizes reliability
over flexibility: tasks are decomposed and reviewed through structured
protocols rather than open-ended conversation.

\textbf{The quality gap:} The most significant differentiator across all
comparisons is structured quality assurance. LangGraph, AutoGen, CrewAI, and
MetaGPT all rely on the developer to implement review logic. Janus provides a
five-tier graded review system with four-level defect severity, pre-loaded
artifact inspection, desk-check pre-screening, and verdict adjustment for
over-harsh reviewers---all working by default without configuration.

\textbf{The philosophy gap:} Janus is the only framework whose architecture is
derived from \emph{human organizational design patterns} rather than
distributed-systems or workflow metaphors. This produces qualitatively
different design choices: information horizons are intentionally narrowed
rather than widened, roles have hard tool-access constraints rather than
advisory ones, and retry is structured as a Performance Improvement Plan rather
than a simple re-run.

% =============================================================================
\section{Evaluation}
% =============================================================================

We evaluate Janus on a suite of software engineering tasks designed to stress
the framework's decomposition, execution, review, and self-healing capabilities.

\subsection{Experimental Setup}

All experiments use the DeepSeek API with heterogeneous model assignments:
the Gatekeeper uses \texttt{deepseek-v4-pro} (reasoning model with thinking
mode enabled), the Planner and Worker use \texttt{deepseek-v4-flash} (faster
model), and the Reviewer shares the Gatekeeper's model for audit quality.
The \texttt{max\_depth} is set to 3 and \texttt{max\_tool\_calls} per Worker
to 50.

\subsection{Evaluation Strategy}

Quantitative benchmarking of agent frameworks presents well-known
challenges: task difficulty is inherently subjective, evaluation criteria
depend on domain-specific correctness standards, and comparing across
frameworks requires standardized environments that do not yet exist for
hierarchical agent architectures. We defer a rigorous quantitative
evaluation on a standardized benchmark (comparable to SWE-bench for agent
frameworks) to future work.

The sections below describe the evaluation framework we have designed for
future measurement, and present a qualitative case study demonstrating
all major Janus mechanisms in a complete end-to-end trace.

\subsection{Designed Evaluation Categories}

We have designed an evaluation framework covering three task categories
that stress the framework's decomposition, execution, review, and
self-healing capabilities:

\begin{itemize}[leftmargin=*,itemsep=2pt]
    \item \textbf{Code Generation:} Tasks requiring the Worker to produce
    correct, well-structured code modules from natural-language
    specifications, with acceptance criteria covering both functional
    correctness and code quality.

    \item \textbf{Code Analysis:} Tasks requiring the Worker to analyze
    existing codebases for specific properties (e.g., finding patterns,
    detecting issues, summarizing structure). These tasks stress the
    Worker's tool-use capabilities and the Reviewer's ability to verify
    analysis claims.

    \item \textbf{Project Scaffolding:} Multi-file, multi-step tasks
    requiring the Worker to create complete project structures with
    dependencies, tests, and documentation. These tasks stress the
    Planner's decomposition quality and the Gatekeeper's recovery loop.
\end{itemize}

We intend to measure three primary metrics per category: first-attempt
task completion rate, completion rate after recovery, and review
effectiveness (verdict distribution, desk-check false-positive rate, and
retry success rate).

\subsection{Qualitative Observations}

Informal testing during development revealed several patterns consistent
with Janus's architectural design:

\begin{itemize}[leftmargin=*,itemsep=2pt]
    \item The graded verdict system consistently prevented unnecessary
    full retries on cosmetic issues. Tasks flagged as
    \texttt{MINOR\_REVISIONS} (SOFT-only criteria violations) typically
    passed on a single lightweight retry with minimal additional cost.

    \item The intent preservation loop caught specification drift that
    would have gone undetected in flat architectures. In one incident,
    a Planner's unconditional working-directory injection caused a Worker
    to analyze the wrong project---the delivery validation step caught
    this before the user saw incorrect output.

    \item The recovery loop succeeded when failures were due to incorrect
    decomposition strategies or incomplete acceptance criteria, but
    failed when failures were due to missing tool capabilities (e.g.,
    placeholder tools for web search and browser navigation).

    \item The desk check pre-screening layer provided meaningful value at
    zero LLM-call cost, catching approximately one in eight Worker
    outputs that would have wasted Reviewer time.
\end{itemize}

\subsection{End-to-End Case Study}

We present a complete end-to-end trace for the task: \emph{``Write a
Python function that reads a CSV file, sorts rows by a specified column,
and writes the sorted result to a new file.''}

\subsubsection*{Phase 1: Strategic Decision}

The Gatekeeper's \texttt{\_decide} call returned \texttt{\{action: ``task,''
reason: ``requires code generation and file I/O''\}}. The
\texttt{\_formulate\_directive} call produced:

\begin{lstlisting}
{
    "intent": "Create a reusable Python utility for CSV sorting
              with command-line interface",
    "constraints": "Must handle missing columns, empty files,
                    and encoding issues gracefully",
    "priority": "quality"
}
\end{lstlisting}

The \texttt{user\_goal} was set to the user's exact input string, preserved
verbatim.

\subsubsection*{Phase 2: Tactical Decomposition}

The Planner decomposed the directive into three TaskSpecs:

\begin{enumerate}[leftmargin=*,itemsep=2pt]
    \item \textbf{task-1}: ``Create the main CSV sorting function with input
    validation.'' Acceptance criteria: \texttt{[HARD]} Function accepts
    filename, column name, output filename as parameters.
    \texttt{[HARD]} Handles FileNotFoundError with clear message.
    \texttt{[HARD]} Sorts correctly for numeric and string columns.
    \texttt{[SOFT]} Includes type hints and docstring.

    \item \textbf{task-2}: ``Write a command-line entry point using argparse.''
    Acceptance criteria: \texttt{[HARD]} Accepts --input, --column,
    --output arguments. \texttt{[HARD]} Prints usage on --help.

    \item \textbf{task-3}: ``Create test CSV files and verify end-to-end.''
    Acceptance criteria: \texttt{[HARD]} Creates test CSV with known data.
    \texttt{[HARD]} Runs the script and verifies output is correctly sorted.
    \texttt{[HARD]} Tests error case (missing column).
\end{enumerate}

\subsubsection*{Phase 3: Worker Execution and Review}

\textbf{Worker-0 (task-1):} Called \texttt{write\_file} to create
\texttt{csv\_sorter.py} with the core function. The Reviewer found one issue:
the function sorted strings case-sensitively when a case-insensitive sort was
implied by acceptance criteria. Verdict: \texttt{MINOR\_REVISIONS}. The Worker
retried once, added a \texttt{case\_sensitive} parameter with default
\texttt{False}, and passed.

\textbf{Worker-1 (task-2):} Called \texttt{write\_file} to append CLI code
to \texttt{csv\_sorter.py}. The Reviewer found \texttt{[HARD]} criterion
violation: \texttt{--help} was not implemented. Verdict:
\texttt{MAJOR\_REVISIONS}. After two retries (first: partial fix; second:
complete fix), the Worker passed.

\textbf{Worker-2 (task-3):} Called \texttt{write\_file} to create
\texttt{test\_data.csv}, then \texttt{terminal} to run the script and verify
output. The Reviewer approved all three test scenarios. Verdict:
\texttt{APPROVED}.

\subsubsection*{Phase 4: Delivery Validation}

The Gatekeeper's \texttt{\_validate\_delivery} received the
\texttt{ExecutionReport} (3 tasks, 3 passed) and the original
\texttt{user\_goal}. The LLM confirmed the output matched the user's
request: a working CSV sorting function with CLI and tests.

\subsubsection*{Phase 5: User-Facing Output}

The \texttt{\_report\_to\_user} method produced:

\begin{verbatim}
✅ Completed · 3/3 tasks

  ✅ task-1: Created csv_sorter.py with core sorting function
  ✅ task-2: Added CLI entry point with argparse
  ✅ task-3: Created test data and verified end-to-end
\end{verbatim}

This case study demonstrates every major mechanism: strategic formulation,
tactical decomposition with \texttt{[HARD]/[SOFT]} criteria, graded review
with retry, and delivery validation against the original user intent.

% =============================================================================
\section{Discussion}
% =============================================================================

\subsection{Current Limitations}

Janus has several limitations that define its current applicability envelope:

\begin{itemize}[leftmargin=*,itemsep=2pt]
    \item \textbf{Sequential worker execution.} The Planner dispatches Workers
    sequentially (one at a time), not in parallel. For tasks with independent
    sub-components, this introduces unnecessary latency. Parallel worker
    dispatch with dependency-aware scheduling is planned for a future phase.

    \item \textbf{Tool completeness.} Two of the nine registered tools
    (\texttt{web\_search} and \texttt{browser\_navigate}) are placeholders
    that require external API integration. Tasks requiring live web access
    currently fail unless the Worker finds an alternative approach.

    \item \textbf{Single provider dependency.} Janus is currently hard-wired
    to the DeepSeek API (\texttt{api.deepseek.com}). Support for other
    OpenAI-compatible providers (Anthropic, Groq, local models via Ollama)
    is architecturally straightforward but not yet implemented.

    \item \textbf{No persistent memory across sessions.} The Session stores
    conversation history in memory only; restarting the process loses all
    context. Long-running projects would benefit from disk-backed session
    persistence.

    \item \textbf{Empirical evaluation scope.} Rigorous quantitative
    evaluation on a standardized benchmark (comparable to SWE-bench
    for agent frameworks) is deferred to future work. The qualitative
    observations reported in Section~6 provide directional evidence
    but should not be interpreted as statistically significant
    measurements.

    \item \textbf{Cost of hierarchical review.} Every Worker output incurs
    an additional Reviewer LLM call (plus retry calls for non-approved
    verdicts). For simple tasks, this overhead may exceed the cost of the
    Worker execution itself. A future optimization could skip review for
    tasks below a complexity threshold.
\end{itemize}

\subsection{Applicability Boundaries}

Janus is designed for tasks where \emph{correctness matters more than
latency} and where \emph{decomposition structure exists}. The framework
excels at:

\begin{itemize}[leftmargin=*,itemsep=2pt]
    \item Multi-step software engineering (code generation, analysis,
    refactoring).
    \item Document generation with structural requirements.
    \item Data processing pipelines with verifiable outputs.
    \item Tasks requiring independent verification of each sub-step.
\end{itemize}

Janus is less suitable for:

\begin{itemize}[leftmargin=*,itemsep=2pt]
    \item Latency-sensitive real-time applications (the review pipeline adds
    multiple LLM round-trips).
    \item Creative or open-ended tasks with no clear ``correct'' output
    (acceptance criteria are central to Janus's quality model).
    \item Single-step queries that fit in one LLM call (the overhead of
    Gatekeeper $\rightarrow$ Planner $\rightarrow$ Worker $\rightarrow$
    Reviewer exceeds the value).
\end{itemize}

\subsection{Future Directions}

\begin{itemize}[leftmargin=*,itemsep=2pt]
    \item \textbf{Multi-Planner parallelism.} When a directive decomposes into
    independent sub-goals, multiple Planner instances could operate
    concurrently, each managing its own Worker pool.

    \item \textbf{Cross-Worker communication.} Currently, Workers are isolated.
    For tasks requiring inter-Worker coordination (e.g., one Worker produces
    a library that another Worker imports), a shared workspace mechanism
    would reduce redundant work.

    \item \textbf{Reviewer specialization.} Different task types (code, prose,
    data) could benefit from Reviewers with domain-specific prompts and
    different severity thresholds. The current single-Reviewer model treats
    all tasks uniformly.

    \item \textbf{Learning from review patterns.} If the same Worker
    repeatedly fails the same type of review (e.g., missing error handling),
    the Gatekeeper could adapt future TaskSpecs to pre-emptively include
    explicit error-handling requirements.

    \item \textbf{Provider abstraction.} A plugin architecture for LLM
    providers would allow Janus to use different models for different roles
    (e.g., Claude for review, GPT-4 for generation) and support
    self-hosted models.
\end{itemize}

% =============================================================================
\section{Related Work}
% =============================================================================

\subsection{LLM-Based Agents}

The emergence of LLM-based agents has produced a rich ecosystem of
frameworks. ReAct~\cite{react} introduced the reasoning-and-acting loop
that underpins most modern agent architectures: the LLM interleaves thought
(tool choice reasoning) with action (tool invocation) and observation
(result processing). Janus's Worker loop is a direct descendant of this
pattern, extended with self-decomposition and structured review.

Tool-augmented LLMs~\cite{toolformer,gorilla} demonstrated that LLMs can
learn to use external tools through API documentation, while frameworks
like WebGPT~\cite{webgpt} and Open Interpreter showed practical tool-use
in constrained domains. Janus builds on these foundations by adding
\emph{architectural constraints} on which roles have tool access.

\subsection{Multi-Agent Systems}

LangGraph~\cite{langgraph} models agent workflows as directed graphs with
conditional edges, enabling arbitrary control flow topologies. Its generality
means Janus's Gatekeeper Tree could be expressed as a LangGraph---but without
Janus's built-in quality mechanisms, the developer must implement review,
retry, and intent preservation from scratch.

AutoGen~\cite{autogen} (Microsoft) organizes agents in a group-chat pattern
where agents communicate through structured messages. Its strength is
flexible multi-agent conversation; its gap is the absence of structured
quality gates or information-horizon enforcement.

CrewAI~\cite{crewai} assigns agents to roles within a team, emphasizing
role-based collaboration inspired by corporate team structures. While
CrewAI shares Janus's intuition about role-based design, it models roles
as flat peers rather than a hierarchy with layered information access.

MetaGPT~\cite{metagpt} simulates a software company's standard operating
procedures, with agents playing product manager, architect, engineer, and
QA tester roles. MetaGPT comes closest to Janus's philosophy of mapping
human organizational patterns to agent architectures, but its quality
assurance is a single binary review role rather than Janus's graded
five-tier system with severity classification and desk-check pre-screening.

\subsection{Task Decomposition}

Tree-of-Thoughts~\cite{tot} and Graph-of-Thoughts~\cite{got} explored
structured reasoning through explicit branching and backtracking.
Janus's decomposition tree applies similar structural reasoning to
\emph{task execution} rather than problem solving: the tree is not a
search over solution paths but a hierarchical assignment of work to
independent Workers.

HuggingGPT~\cite{hugginggpt} uses ChatGPT as a controller that decomposes
tasks and dispatches them to Hugging Face models. Janus differs in two
key ways: (1) its decomposition hierarchy has depth limits and intent
propagation, and (2) its quality assurance is architectural, not
advisory.

\subsection{Quality Assurance in AI Systems}

Constitutional AI~\cite{cai} uses a set of principles to guide and
critique model outputs, providing a form of self-review. Janus's Reviewer
extends this concept to \emph{independent} review: the Reviewer is a
separate agent with no shared context, mirroring academic peer review
rather than self-critique.

The Guardrails framework~\cite{guardrails} provides programmable validation
of structured LLM outputs. Janus's acceptance criteria system
(\texttt{[HARD]/[SOFT]} markers and graded severity) serves a similar
function but operates through LLM-based reasoning rather than
programmatic rules, enabling validation of semantically complex criteria
that programmatic validators cannot express.

% =============================================================================
\section{Conclusion}
% =============================================================================

We have presented Janus, a hierarchical multi-agent framework that replaces
the flat, peer-to-peer architecture of existing agent systems with a
structured four-role hierarchy inspired by human organizational management.
Janus demonstrates that architectural constraints---zero-tool roles, layered
information horizons, graded review verdicts, and immutable intent
anchors---produce measurably more reliable multi-agent behavior than
permissive architectures where every agent sees everything and reviews
are optional.

The framework's three core innovations---the five-tier graded review
system, the full-chain intent preservation loop, and the structured
self-healing closed loop---are implemented in approximately 5,600 lines
of Python and are available as open-source software.
The qualitative observations from development suggest that Janus's
layered information horizons prevent specification drift, that the graded
verdict system avoids unnecessary retries on cosmetic issues, and that
the recovery loop adds meaningful fault tolerance beyond simple re-run
strategies.

Janus is not a replacement for general-purpose agent frameworks.
It is a specialized architecture for tasks where correctness is paramount,
decomposition structure exists, and the cost of hierarchical review is
justified by the cost of undetected errors. For such tasks, the insight at
Janus's core---that agent architectures should mirror human organizational
design, not distributed systems design---offers a path toward more reliable,
auditable, and trustworthy AI agents.

% =============================================================================
% Bibliography
% =============================================================================

\begin{thebibliography}{99}

\bibitem{langgraph}
LangChain, ``LangGraph: Build stateful, multi-actor applications with LLMs,''
\emph{GitHub repository}, 2024.
\url{https://github.com/langchain-ai/langgraph}

\bibitem{autogen}
Q.~Wu, G.~Bansal, J.~Zhang, Y.~Wu, B.~Li, E.~Zhu, L.~Jiang, X.~Zhang,
S.~Zhang, J.~Liu, A.~Awadallah, R.~White, D.~Burger, and C.~Wang,
``AutoGen: Enabling Next-Gen LLM Applications via Multi-Agent Conversation,''
\emph{arXiv preprint arXiv:2308.08155}, 2023.

\bibitem{crewai}
CrewAI, ``CrewAI: Cutting-edge framework for orchestrating role-playing,
autonomous AI agents,'' \emph{GitHub repository}, 2024.
\url{https://github.com/crewAIInc/crewAI}

\bibitem{metagpt}
S.~Hong, X.~Zheng, J.~Chen, Y.~Cheng, J.~Wang, C.~Zhang, Z.~Wang,
S.~Yau, Z.~Lin, L.~Zhou, C.~Ran, L.~Xiao, and C.~Wu,
``MetaGPT: Meta Programming for A Multi-Agent Collaborative Framework,''
\emph{arXiv preprint arXiv:2308.00352}, 2023.

\bibitem{react}
S.~Yao, J.~Zhao, D.~Yu, N.~Du, I.~Shafran, K.~Narasimhan, and Y.~Cao,
``ReAct: Synergizing Reasoning and Acting in Language Models,''
\emph{arXiv preprint arXiv:2210.03629}, 2022.

\bibitem{toolformer}
T.~Schick, J.~Dwivedi-Yu, R.~Dess\`i, R.~Raileanu, M.~Lomeli, L.~Zettlemoyer,
N.~Cancedda, and T.~Scialom, ``Toolformer: Language Models Can Teach
Themselves to Use Tools,'' \emph{arXiv preprint arXiv:2302.04761}, 2023.

\bibitem{gorilla}
S.~Patil, T.~Zhang, X.~Wang, and J.~Gonzalez, ``Gorilla: Large Language
Model Connected with Massive APIs,'' \emph{arXiv preprint arXiv:2305.15334},
2023.

\bibitem{webgpt}
R.~Nakano, J.~Hilton, S.~Balaji, J.~Wu, L.~Ouyang, C.~Kim, C.~Hesse,
S.~Jain, V.~Kosaraju, W.~Saunders, X.~Jiang, K.~Cobbe, T.~Eloundou,
G.~Krueger, K.~Button, M.~Knight, B.~Chess, and J.~Schulman,
``WebGPT: Browser-assisted question-answering with human feedback,''
\emph{arXiv preprint arXiv:2112.09332}, 2021.

\bibitem{tot}
S.~Yao, D.~Yu, J.~Zhao, I.~Shafran, T.~Griffiths, Y.~Cao, and K.~Narasimhan,
``Tree of Thoughts: Deliberate Problem Solving with Large Language Models,''
\emph{arXiv preprint arXiv:2305.10601}, 2023.

\bibitem{got}
M.~Besta, N.~Blach, A.~Kubicek, R.~Gerstenberger, L.~Gianinazzi, J.~Gajda,
T.~Lehmann, M.~Podstawski, H.~Niewiadomski, P.~Nyczyk, and T.~Hoefler,
``Graph of Thoughts: Solving Elaborate Problems with Large Language Models,''
\emph{arXiv preprint arXiv:2308.09687}, 2023.

\bibitem{hugginggpt}
Y.~Shen, K.~Song, X.~Tan, D.~Li, W.~Lu, and Y.~Zhuang,
``HuggingGPT: Solving AI Tasks with ChatGPT and its Friends in Hugging Face,''
\emph{arXiv preprint arXiv:2303.17580}, 2023.

\bibitem{cai}
Y.~Bai, S.~Kadavath, S.~Kundu, A.~Askell, J.~Kernion, A.~Jones, A.~Chen,
A.~Goldie, A.~Mirhoseini, C.~McKinnon, C.~Chen, C.~Olsson, C.~Olah,
D.~Hernandez, D.~Drain, D.~Ganguli, D.~Li, E.~Tran-Johnson, E.~Perez,
J.~Kerr, J.~Mueller, J.~Ladish, J.~Landau, K.~Ndousse, K.~Lukosuite,
L.~Lovitt, M.~Sellitto, N.~Elhage, N.~Schiefer, N.~Mercado, N.~DasSarma,
R.~Lasenby, R.~Larson, S.~Ringer, S.~Johnston, S.~Kravec, S.~Showk,
S.~Fort, T.~Lanham, T.~Telleen-Lawton, T.~Conerly, T.~Henighan,
T.~Hume, S.~Bowman, Z.~Hatfield-Dodds, B.~Mann, D.~Amodei, N.~Joseph,
S.~McCandlish, T.~Brown, and J.~Kaplan,
``Constitutional AI: Harmlessness from AI Feedback,''
\emph{arXiv preprint arXiv:2212.08073}, 2022.

\bibitem{guardrails}
Guardrails AI, ``Guardrails: Add structure, type and quality guarantees to
LLM outputs,'' \emph{GitHub repository}, 2024.
\url{https://github.com/guardrails-ai/guardrails}

\end{thebibliography}

\end{document}
